\documentclass[tikz, border=6pt]{standalone}
% Shared preamble of the standalone TikZ figures (compiled by scripts/build_figures.py).
\usepackage{amsmath, amssymb}
\usepackage{xeCJK}
\setCJKmainfont[BoldFont={Hiragino Sans W6}]{Hiragino Sans W4}
\usetikzlibrary{arrows.meta, positioning, decorations.pathreplacing, calc, patterns}
\definecolor{rod}{RGB}{70,110,170}
\definecolor{rodb}{RGB}{200,120,50}
\definecolor{okc}{RGB}{60,140,80}
\definecolor{ngc}{RGB}{190,70,60}
\definecolor{boxc}{RGB}{235,240,248}
\tikzset{
  >={Stealth[length=2.4mm]},
  every node/.style={font=\small},
  box/.style={draw, rounded corners=2pt, fill=boxc, align=flush center, inner sep=4pt},
  rodA/.style={draw, fill=rod!35},
  rodB/.style={draw, fill=rodb!35},
  note/.style={font=\footnotesize, align=flush center},
  expl/.style={font=\small, align=flush left},
  brace/.style={decorate, decoration={brace, amplitude=4pt}},
}

\begin{document}
\begin{tikzpicture}
  % (a) comparison
  \node[anchor=west, font=\bfseries] at (-0.6, 1.5) {(a) 比較 $d \preceq e$};
  \draw[dashed, gray] (0,-1.0) -- (0,1.0);
  \node[note, gray, anchor=north] at (0,-1.0) {同時に始める};
  \draw[rodA] (0, 0.35) rectangle (2.6, 0.75) node[midway] {$d$：砂時計};
  \draw[rodB] (0,-0.55) rectangle (4.0,-0.15) node[midway] {$e$：振り子 10 回};
  \draw[dashed, gray] (2.6, 0.35) -- (2.6,-0.75);
  \node[note, anchor=west] at (2.7, 0.55) {$d$ が先に終わる};

  % (b) concatenation
  \begin{scope}[xshift=7.6cm]
    \node[anchor=west, font=\bfseries] at (-0.6, 1.5) {(b) 連結 $d \circ e$};
    \draw[rodA] (0, 0) rectangle (2.6, 0.4) node[midway] {$d$};
    \draw[rodB] (2.6, 0) rectangle (6.6, 0.4) node[midway] {$e$};
    \node[note, above] at (2.6, 0.45) {$d$ が終わったら $e$ を始める};
    \draw[brace, decoration={mirror}] (0,-0.1) -- (6.6,-0.1) node[midway, below=5pt] {$\mathrm{wait}_{d \circ e} = \mathrm{wait}_e \circ \mathrm{wait}_d$};
  \end{scope}

  \node[expl, anchor=west, text width=14cm] at (-0.6,-1.9) {比較と連結が長さと同じ法則を満たせば、長さの定理がそのまま時間の目盛りを与える。同じ比較の結果を与える二つの時計の読みは比例する。};
\end{tikzpicture}
\end{document}
